\documentclass[11pt]{article}

\usepackage[utf8]{inputenc}
\usepackage[T1]{fontenc}

\usepackage[slovene]{babel}
\usepackage{lmodern}
\usepackage{amsmath}
\usepackage{amsfonts}
\usepackage{amsthm}
\usepackage{mathtools}

\begin{document}

V to datoteko vnesite svoje odgovore na vprašanja iz datoteke \texttt{kviz.pdf}
Pri nekaterih nalogah je možnih več pravilnih odgovorov. Izberite vse.

\begin{enumerate}
    \item Pravilni odgovori so: % tu napišite vse pravilne odgovore
    
    \item Pravilni odgovori so: % tu napišite vse pravilne odgovore
    
    \item Pravilni odgovori so:  % tu naštejte vse pravilne odgovore
    % (a)
      \[ e^{i \pi} + 1 = 0 \]
    % (b)
      \begin{center}
          \( e^{i \pi} + 1 = 0 \)
      \end{center}
    % (c)
      \begin{equation*}
          e^{i \pi} + 1 = 0
      \end{equation*}

    \item Pravilni odgovori so: % tu naštejte vse pravilne odgovore
        \begin{enumerate}
          % (a)
            \item {
Lorem ipsum dolor sit amet, consectetur adipiscing 
elit. Ut id viverra ligula. Phasellus vehicula lorem 
vitae luctus dignissim. 
                
Sed ac justo commodo, fringilla urna ac, efficitur 
leo. Praesent dui odio, accumsan ac sapien nec, 
interdum volutpat est. 
            }
          % (b)
            \item {
Lorem ipsum dolor sit amet, consectetur adipiscing 
elit. Ut id viverra ligula. Phasellus vehicula lorem 
vitae luctus dignissim. \\
Sed ac justo commodo, fringilla urna ac, efficitur 
leo. Praesent dui odio, accumsan ac sapien nec, 
interdum volutpat est.             }
          % (c)
            \item {
Lorem ipsum dolor sit amet, consectetur adipiscing 
elit. Ut id viverra ligula. Phasellus vehicula lorem 
vitae luctus dignissim. \par
Sed ac justo commodo, fringilla urna ac, efficitur 
leo. Praesent dui odio, accumsan ac sapien nec, 
interdum volutpat est.             }
        \end{enumerate}
    
      \item Pravilni odgovori so: % tu naštejte vse pravilne odgovore
        \begin{enumerate}
        % (a)
          \item "Primer 1"
        % (b)
          \item ``Primer 2''
        % (c)
          \item ">Primer 3"<
        % (d)
          \item "`Primer 4"'
        \end{enumerate}

    \item Pravilni odgovori so: % tu naštejte vse pravilne odgovore
    \begin{enumerate}
        % (a)
        \item \LaTeX je dokumentni sistem.
        % (b)
        \item {\LaTeX} je dokumentni sistem.
        % (c)
        \item \LaTeX{} je dokumentni sistem.
    \end{enumerate}

    \item Pravilni odgovori so: % tu naštejte vse pravilne odgovore
    \begin{enumerate}
        % (a)
        \item prof. Novak
        % (b)
        \item prof.~Novak
        % (c)
        \item prof.\ Novak
    \end{enumerate}

    \item Pravilni odgovori so: % tu naštejte vse pravilne odgovore
    \begin{enumerate}
        % (a)
        \item
Za kateti in hipotenuzo velja
\[
  a^2 + b^2 = c^2,
\]
kjer je \(c\) hipotenuza.
        % (b)
        \item
Za kateti in hipotenuzo velja
\[
  a^2 + b^2 = c^2,
\]

kjer je \(c\) hipotenuza.
        % (c)
        \item
Za kateti in hipotenuzo velja

\[
  a^2 + b^2 = c^2,
\]
kjer je \(c\) hipotenuza.
    \end{enumerate}

    \item Pravilni odgovori so: % tu naštejte vse pravilne odgovore
    \begin{enumerate}
        % (a)
        \item Vsaka surjekcija ima prerez\footnote{Opomba.}, zato je svet lepši.
        % (b)
        \item Vsaka surjekcija ima prerez,\footnote{Opomba.} zato je svet lepši.
    \end{enumerate}

    \item Pravilni odgovori so: % tu naštejte vse pravilne odgovore
    \begin{enumerate}
        % (a)
        \item 50 %
        % (b)
        \item 50 \%
        % (c)
        \item 50~\%
    \end{enumerate}

    \item Pravilni odgovori so: % tu naštejte vse pravilne odgovore
    \begin{enumerate}
        % (a)
        \item Naj bo x realno število.
        % (b)
        \item Velja \(x\) + \(y\) = \(y\) + \(x\).
        % (c)
        \item Velja \(x + y = y + x\).
    \end{enumerate}

    \item Pravilni odgovori so: % tu naštejte vse pravilne odgovore
    \begin{enumerate}
        % (a)
        \item \(x_1\), \dots, \(x_n\)
        % (b)
        \item \(x_1, \ldots, x_n\)
    \end{enumerate}

    \item Pravilni odgovori so: % tu naštejte vse pravilne odgovore
    \begin{enumerate}
        % (a)
        \item
Velja
\[
  a^2 + b^2 = c^2.
\]
        % (b)
        \item
Velja
\[
  a^2 + b^2 = c^2
\]
.
        % (c)
        \item
Velja
\[
  a^2 + b^2 = c^2
\]
    \end{enumerate}

    \item Pravilni odgovori so: % tu naštejte vse pravilne odgovore
    \begin{enumerate}
        % (a)
        \item \(<x, y>\)
        % (b)
        \item \(\langle x, y \rangle\)
    \end{enumerate}

    \item Pravilni odgovori so: % tu naštejte vse pravilne odgovore
    \begin{enumerate}
        % (a)
        \item \[ (\frac{a}{b})^n \]
        % (b)
        \item \[ \left(\frac{a}{b}\right)^n \]
        % (c)
        \item \[ \Big(\frac{a}{b}\Big)^n \]
    \end{enumerate}

    \item Pravilni odgovori so: % tu naštejte vse pravilne odgovore
    \begin{enumerate}
        % (a)
        \item Prva vrstica.
\begin{align*}
  (x + y)^2 &= (x + y)(x + y) \\
  &= x^2 + 2 x y + y^2 \\
\end{align*}
        Vrstica za enačbama.
        % (b)
        \item Prva vrstica.
\begin{align*}
  (x + y)^2 &= (x + y)(x + y) \\
  &= x^2 + 2 x y + y^2
\end{align*}
        Vrstica za enačbama.
    \end{enumerate}

    \item Pravilni odgovori so: % tu naštejte vse pravilne odgovore
    \begin{enumerate}
        % (a)
        \item Prva vrstica.
\[ a + b = b + a, \]
\[ a b = b a, \]
\[ a (b + c) = a b + a c. \]
        % (b)
        \item Prva vrstica.
\begin{gather*}
  a + b = b + a, \\
  a b = b a, \\
  a (b + c) = a b + a c.
\end{gather*}
    \end{enumerate}

    \item Pravilni odgovori so: % tu naštejte vse pravilne odgovore

\end{enumerate}

\end{document}
